% Concise revised introduction for Overleaf.
% Citations are written as author-year placeholders for later integration.

\documentclass[11pt]{article}
\usepackage[margin=1in]{geometry}
\usepackage{amsmath}
\usepackage[version=4]{mhchem}
\usepackage[T1]{fontenc}
\usepackage{lmodern}

\title{Precession-Phase Organization of DO-Like Asian Monsoon Transitions}
\author{}
\date{}

\begin{document}
\maketitle

\section*{Introduction}

Millennial-scale climate variability (MCV) is a prominent mode of abrupt climate change in glacial climate systems. Its best-known expression is the Dansgaard--Oeschger (DO) sequence in Greenland ice cores, where rapid stadial--interstadial transitions reflect major reorganizations of the North Atlantic climate system (Rasmussen et al., 2014). However, MCV was not restricted to Greenland or to the last glacial cycle. Marine, ice-core, and speleothem archives suggest that abrupt millennial variability persisted across multiple Pleistocene glacial cycles and was modulated by changing climate background states (Barker et al., 2011; Sun et al., 2021; Hodell et al., 2023). Speleothems are particularly valuable in this context because their U--Th chronologies allow abrupt hydroclimate changes to be compared with ice-core and orbital timescales. During the last glacial period, independently dated speleothem records show near-synchronous abrupt hydroclimate changes in monsoon and mid-latitude regions during Greenland stadial--interstadial transitions (Corrick et al., 2020). Chinese speleothem records therefore provide a precisely dated archive of the Asian monsoon expression of DO-like MCV (Wang et al., 2001; Wang et al., 2008).

Chinese speleothem \(\delta^{18}\mathrm{O}\) records also contain one of the clearest orbital-scale monsoon signals of the late Quaternary. Multi-cave composites from Hulu, Dongge, Sanbao, and related sites show dominant precession-band variability that is closely linked to boreal summer insolation, with millennial-scale monsoon events superimposed on this orbitally paced background (Wang et al., 2008; Cheng et al., 2016). The 640 kyr Chinese speleothem stack extended this framework across several glacial cycles and suggested that insolation partly sets the pace of millennial-scale monsoon variability (Cheng et al., 2016). Yet the interpretation of cave \(\delta^{18}\mathrm{O}\) remains non-trivial: it is not a simple local rainfall proxy, but reflects large-scale monsoon circulation, moisture-source changes, upstream rainout, and regional hydroclimate dynamics (Liu et al., 2014; Cheng et al., 2021; Wang et al., 2022). Thus, the precession signal in Chinese cave records is best interpreted as a large-scale Asian monsoon and water-cycle signal rather than as local precipitation alone.

A key event-based question nevertheless remains unresolved. Previous work has mainly examined continuous \(\delta^{18}\mathrm{O}\) variability through spectra, phase relationships, amplitude modulation, or correlations with insolation. These analyses show that the monsoon background state is strongly paced by precession, but they do not establish whether objectively identified abrupt transition timings are themselves organized by orbital phase. Abrupt transitions could occur preferentially within certain orbital windows, but they could also arise from internally generated threshold dynamics whose amplitudes are later modulated by orbital-scale background changes (Ditlevsen and Ditlevsen, 2005; Lohmann and Ditlevsen, 2018). Recent model results further suggest that astronomical configuration may directly influence DO-like oscillations through precession- and obliquity-related changes in tropical hydrology, sea ice, and ocean circulation (Zhang et al., 2021). These results motivate a proxy-based test of whether speleothem-recorded abrupt Asian monsoon transitions cluster at preferred orbital phases, and whether any phase dependence is independent of slow ice-volume and greenhouse-gas background states.

Here we test this hypothesis using the 0--640 ka sequence of weak and strong Asian monsoon starts identified by Rousseau et al. (2023) from the Chinese speleothem composite of Cheng et al. (2016). We treat these events as DO-like Asian monsoon transitions recorded in speleothems, rather than as Greenland-defined DO events. We first evaluate whether the transition sequences are consistent with stationary random occurrence. We then use circular Rayleigh tests to assess clustering in precession phase, and binned Poisson hazard generalized linear models to test whether precession phase adds predictive information beyond LR04 ice volume and atmospheric \(\ce{CO2}\). This framework separates the well-established precession pacing of continuous speleothem \(\delta^{18}\mathrm{O}\) variability from the more specific question of whether abrupt transition occurrence is itself phase-organized.

We find that weak and strong monsoon starts are inconsistent with simple stationary randomness and show significant, but opposite, precession-phase preferences. Poisson hazard models further indicate that precession phase remains informative after accounting for LR04 and \(\ce{CO2}\), whereas obliquity phase and other candidate forcings do not provide comparable independent contributions. These results suggest that DO-like Asian monsoon transitions were not merely embedded within a precession-paced monsoon background; their occurrence probability was itself organized by precession phase. We therefore interpret precession as a phase-dependent boundary condition that modulates transition likelihood together with slower ice-volume and greenhouse-gas background states, while preserving the distinction between Greenland-defined DO events and DO-like monsoon transitions recorded in Chinese speleothems.

\end{document}
